\section{What determines whether a sound evokes orienting, vigilance, freezing, flight, or approach?}
\label{sec:q4}

% Part II — Assessment and behavioral response (Q4).
% See docs/OUTLINE.md — Q4.
% Draft 1. Style models: Magrath + Blumstein (agents/STYLE.md).

The ecological consequence of a sound depends on the behavior it elicits, and the behaviors available to a rodent differ in latency, cost, and what they imply about assessment. Acoustic startle is a short-latency reflex; orienting and vigilance gather information at the cost of foraging time; freezing reduces detection by a predator but forgoes escape; flight to a refuge interrupts foraging and may expose the animal; and approach indicates that a sound is attractive rather than aversive. These responses are not interchangeable measures of a single fear state. Startle in particular habituates according to the timing of stimulation \citep{Davis1970} and shows that a sound was detected, not that a threat was assessed or that any defense will be sustained. The question for this section is which properties of the sound, the receiver, and its situation determine which response is expressed.

Threat-sensitive predator avoidance predicts that response intensity should scale with the risk a sound indicates \citep{Helfman1989}, and sciurid field studies support this prediction with graded acoustic variation. Marmots respond more to calls given at higher rates \citep{BlumsteinArmitage1997} and to calls with added nonlinear components \citep{BlumsteinRecapet2009}.
% VERIFY: BlumsteinArmitage1997 and BlumsteinRecapet2009 (not full-text checked): response scaling with call rate and nonlinearity.
In Richardson's ground squirrels, playback evokes a range of responses, from changes in posture and vigilance duration to escape and burrow entry, that depend on the caller and call structure \citep{Hare1998}. Graded intensity is not the only form of variation. In Gunnison's prairie dogs, playback of alarm calls given to different predators reproduced qualitatively different escape actions: diving into the burrow for calls given to humans, running to the burrow and standing at its rim for calls given to coyotes, and adopting an alert posture in place for calls given to dogs \citep{KiriazisSlobodchikoff2006}. Receivers in this system select among actions rather than adjusting the strength of a single response.

The same sound can also produce different actions depending on the receiver's situation. In laboratory mice, trains of a synthetic 17--20~kHz frequency sweep elicited flight or freezing depending on whether the mouse was in its home cage or an unfamiliar cage and whether it had previously received foot shock \citep{MongeauEtAl2003}. Because a single stimulus was used, the study isolates context and experience from acoustic structure.
% VERIFY: MongeauEtAl2003 (not full-text checked): context- and shock-dependent shift between flight and freezing.
% VERIFY: EilamEtAl1999 and HernandezEtAl2021 (not full-text checked): species and habitat dependence of responses to owl calls.
Across species, owl calls led voles to freeze or flee but changed the behavior of spiny mice little while raising their cortisol \citep{EilamEtAl1999}, and free-living rodents adjusted feeding, movement, and vigilance according to owl species, moonlight, and cover \citep{HernandezEtAl2021}. These results fit the economics of escape (\secref{sec:intro-framework}), in which the action chosen depends on the costs of fleeing and of staying \citep{YdenbergDill1986}. The fit is qualitative; none of these acoustic studies measured the costs that the model requires.

Field studies supply those costs. Alarm-call playback to eastern chipmunks increased vigilance, shortened the run to cover, delayed re-emergence, and reduced food intake \citep{BaackSwitzer2000}. In eastern gray squirrels, predator-call playback during foraging increased vigilance and suppressed and delayed foraging, but did not change the distance squirrels kept from refuge \citep{JayneEtAl2015}.
% VERIFY: BaackSwitzer2000 and JayneEtAl2015 (not full-text checked): foraging costs; distance to refuge unchanged across acoustic conditions.
The second result shows that receivers can pay a foraging cost without adjusting their spatial exposure, so a single endpoint can miss the response that a sound actually produced.
% VERIFY: KittendorfDantzer2021 (not full-text checked): playbacks paired with an approaching human.
Studies that combine acoustic playback with an approaching human \citep{KittendorfDantzer2021} cannot isolate the contribution of sound and are not considered further.

Physiological and laboratory measures show that a sound can be processed without the behavior an experimenter expects.
% VERIFY: Mateo2010 (not full-text checked): cortisol increase to trills; heart-rate slowing to whistles.
Belding's ground squirrels mount different physiological responses to alarm calls associated with terrestrial and aerial predators \citep{Mateo2010} (\secref{sec:q2}). In rats, 50-kHz calls elicit approach \citep{WilluhnEtAl2014,OlszynskiEtAl2020} whereas 22-kHz calls slow locomotion and heart rate \citep{OlszynskiEtAl2020}, and naïve rats discriminate 22-kHz calls neurally although they freeze to them only after aversive experience (\secref{sec:q2}, \secref{sec:q3}). A study that scores only freezing would classify naïve rats as unresponsive to 22-kHz calls and spiny mice as unresponsive to owl calls, and both conclusions would be wrong. Assessing whether a sound has been interpreted as a threat therefore requires a set of outcomes: short-latency orienting or startle, locomotion, approach or avoidance, vigilance, autonomic and endocrine measures, and the latency to resume prior activity.

Laboratory responses are valid proxies for ecological decisions only to the extent that the laboratory imposes the trade-offs that shape decisions in the field. Most laboratory assays remove the foraging opportunities, refuges, and social context that determine the cost of responding, and they rarely measure recovery latency, the variable that most directly links a sound to lost foraging. Two gaps follow. We found no study that measured flight initiation distance with acoustic cues alone, and the closest test, in gray squirrels, found no change in refuge distance \citep{JayneEtAl2015}. Nor did we find a systematic rodent test of how hunger, body condition, and other state variables, which the economics of escape treats as central, alter acoustic responses.

\tabref{tab:summary-q4} consolidates these claims for later condensation or figure drafting.

% Intermediate summary matrix for Q4. May later be condensed, moved to an appendix, or redrawn as a figure.
\begin{table}[tbp]
\centering
\footnotesize
\singlespacing
\caption{Working summary for Q4: behavioral outcomes and what each shows. Responses are not interchangeable measures of a single fear state.}
\label{tab:summary-q4}
\begin{tabular}{@{}p{1.15in}p{3.5in}p{1.55in}@{}}
\toprule
\textbf{Outcome} & \textbf{What it indicates} & \textbf{Main determinants} \\
\midrule
Startle & Detection; short-latency reflex & Timing of stimulation \\
\addlinespace[2pt]
Orienting / vigilance & Assessment; costs foraging time & Call rate, caller, risk cues \\
\addlinespace[2pt]
Freezing & Defense without escape & Experience, species, context \\
\addlinespace[2pt]
Flight / burrow entry & Escape selected among options & Predator-specific call category; refuge \\
\addlinespace[2pt]
Approach & Appetitive meaning & Call valence (e.g., 50-kHz) \\
\addlinespace[2pt]
Autonomic / endocrine & Assessment without required overt behavior & Call type, predator class, species \\
\addlinespace[2pt]
Recovery latency / intake & Ecological cost of responding & Field costs absent in most lab assays \\
\bottomrule
\end{tabular}
\end{table}

